\documentclass[11pt]{article}
\usepackage[margin=1in]{geometry}
\usepackage{amsmath,amssymb,amsthm}
\usepackage[T1]{fontenc}
\usepackage{lmodern}
\title{An Internal Contradiction in Conjecture 00000004383}
\author{earthking11}
\date{October 6, 2026}
\begin{document}
\maketitle

\section*{The two assertions}
The conjecture states that the first five nonzero low-degree homotopy groups of the topological Hochschild homology of the sphere spectrum are direct sums of copies of $\mathbb F_p$. It also states that the first nonzero group is cyclic of order four. The first nonzero group is one of the first five, so both descriptions apply to the same group.

\section*{Why they cannot both hold}
The additive group of a field of characteristic $p$ is annihilated by multiplication by $p$. The same is true coordinatewise for any direct sum of copies of that additive group. A cyclic group of order four has a generator $g$ whose order is four, so $p g=0$ implies $4\mid p$.

This contradicts primality. Indeed, write $p=4k=2(2k)$. If $p$ is prime, its factorization property implies either $2=1$ or $2k=1$. The first is false, and the second is impossible because $2k$ is even. Equivalently, no prime is divisible by four. Thus the two descriptions in the conjecture cannot refer to the same group.

\section*{Formalization scope}
The accompanying Lean development proves the direct-sum annihilation bridge, formalizes a cyclic order-four generator by its annihilator criterion, and proves the elementary prime-divisibility contradiction without assuming it. It establishes an internal inconsistency in the stated properties; it does not compute the actual homotopy groups or decide their degrees.

\end{document}
